\subsection{Purpose}
The \tB{back-propagation} algorithm allows the information from the cost to then flow
backwards through the network \tB{in order to compute the gradient}.\\
Actually it refers only to the \tB{method for computing the gradient}, while
\uB{another algorithm such as stochastic gradient descent is used to perform learning}
using this gradient.

\subsection{Modeling}
\paragraph{Chain rule of calculus}
Let $(\bm{x}, \bm{y}) \in \mathbb{R}^{m} \times \mathbb{R}^{n}$ and $z = f(\bm{y})$ 
with $f: \mathbb{R}^{n} \longrightarrow \mathbb{R}$
we have $\nabla_{\bm{x}}z = 
\begin{bmatrix}
    \frac{dz}{dx_{1}}\\
    \vdots \\
    \frac{dz}{dx_{m}}
\end{bmatrix}$,
$\nabla_{\bm{y}}z =
\begin{bmatrix}
    \frac{dz}{dy_{1}}\\
    \vdots \\
    \frac{dz}{dy_{n}}
\end{bmatrix}$ and 
$ \dfrac{\partial\bm{y}}{\partial\bm{x}} = 
\begin{bmatrix}
    \frac{dy_{1}}{x_{1}}  & \cdots & \frac{dy_{1}}{x_{m}}\\
    \vdots                & \ddots & \vdots\\            
    \frac{dy_{n}}{x_{1}}  & \cdots & \frac{dy_{n}}{x_{m}}
\end{bmatrix}$
we have :
\begin{center}
    $\nabla_{\bm{x}}z = \left(\dfrac{\partial\bm{y}}{\partial\bm{x}}\right)^{T}
    \nabla_{\bm{y}}z$
\end{center}

\paragraph{Back-propagation in fully connected MLP}
Consider a network of depth $l$, and the matrices of the model 
$\left(\bm{W}_{j}\right)_{1\leq j\leq l}$ and the bias parameters 
$\left(\bm{b}_{j}\right)_{1\leq j\leq l}$ then $\bm{x}$ the input vector and 
$\bm{y}$ the target vector.

\subparagraph{Forward propagation}
Consider a $l-1$ hidden layer model, assume that we use an unique activation function $f$:
\begin{enumerate}
    \item $\bm{h}_{0} = \bm{x}$ \textit{consider the input as the first layer}
    \item for $k \in \inter{1}{l}$ 
        \begin{itemize}
            \item $\bm{a}_{k} = \bm{b}_{k} + \bm{W}_{k}\bm{h}_{k-1}$  \textit{vector containing all 
                the resulting weighted sum and bias of inputs $k$}
            \item $\bm{h}_{k} = f(\bm{a}_{k})$ \textit{vector containing all the neuron outputs of 
                the considered layer by applying the activation function}
        \end{itemize}
    \item $\hat{\bm{y}} = \bm{h}_{l}$ \textit{output layer}
    \item $J = L(\hat{\bm{y}}, \bm{y}) + \lambda\Omega(\theta)$ \textit{finally compute loss 
        function with a potential regulation term}
        
\end{enumerate}

\subparagraph{Backward}
\begin{enumerate}
    \item $\bm{g} \leftarrow \nabla_{\hat{\bm{y}}}J = \nabla_{\hat{\bm{y}}}L(
        \hat{\bm{y}}, \bm{y})$ \textit{initial gradient vector}
    \item for $k\in\inter{l-1}{1}$
        \begin{itemize}
            \item $\bm{g} \leftarrow \nabla_{\bm{a}_{k}}\bm{J} = \bm{g}\odot f'(\bm{a}_{k})$
                \textit{compute derivative with respect to $\bm{a}_{k}$, Hadamard product is used to imitate the rule: $\left(h\circ f\right)(x)' = h'\left(f(x)\right)f'(x)$}
            \item $\nabla_{\bm{b}_{k}}J = \bm{g} + 
                \lambda\nabla_{\bm{b}_{k}}\Omega(\bm{\theta})$
            \item $\nabla_{\bm{W}_{k}}J = \bm{g}\bm{h}^{T}_{k-1}
                + \lambda\nabla_{\bm{W}_{k}}\Omega(\bm{\theta})$ # TO COMPLETE
            \item $\bm{g} \leftarrow \nabla_{\bm{h}_{k-1}}J = \bm{W}_{k}^{T}\bm{g}$
        \end{itemize}
        
\end{enumerate}


